\section{Related literature}\label{sec:literature}

Three threads of empirical literature converge on the questions this paper addresses: sovereign-credit fundamentals, panel-data methods for causal inference, and the cross-sectional structure of emerging-market returns. Each thread brings its own data, methods and biases. The contribution of the present work is at the intersection, so we sketch all three.

\subsection{Sovereign credit fundamentals}

The empirical sovereign-credit literature dates the modern formulation to \textcite{cantorpacker1996}, who show that ratings are explained by a parsimonious list of macro variables: per-capita income, GDP growth, inflation, external debt, fiscal balance and a sovereign default history dummy. Their cross-sectional regression on a sample of advanced and emerging economies established the headline determinants the field has worked with ever since.

\textcite{hilschernosbusch2010} update that framework with country fundamentals (the volatility of inflation, terms-of-trade movements) and macroeconomic risks, and find that fundamentals price the sovereign spread, not just the rating. The two are not redundant: ratings are scale variables that compress fundamentals into 22 discrete buckets, while spreads embed both fundamentals and a time-varying risk premium that responds to global appetite for emerging-market debt.

\textcite{augustinetal2014} survey the credit-default-swap-sovereign nexus. The main empirical finding from that survey, replicated in subsequent work, is that sovereign CDS spreads load most heavily on global risk-aversion factors at the daily frequency and on country fundamentals at lower frequencies. For our purposes the practical implication is that the slow-moving information in development indicators should show up in ratings (lower-frequency) before it shows up in spreads, and certainly before it shows up in monthly equity returns.

\textcite{kaufmannkraaymastruzzi2010} document that the six Worldwide Governance Indicators --- political stability, rule of law, control of corruption, government effectiveness, voice and accountability, regulatory quality --- explain a substantial share of cross-sectional sovereign-spread variation, even after conditioning on the conventional macro panel. The KPI block in our composite includes all six WGI indicators in the security-stability sector, so part of what we measure has been measured before. The novelty is in the joint treatment of governance \emph{together with} the five other sectors (education, energy, R\&I, health, housing) and the principal-component aggregation that gives each sector a single, interpretable score.

\textcite{reinhart_rogoff2009} and \textcite{reinhart_rogoff2010} document the recurrence of sovereign-default cycles and underline that financial stress predictably follows macroeconomic build-ups --- in particular, that the run-up of public-debt-to-GDP above approximately 90~\% is associated with substantially lower subsequent growth in a large cross-country panel. Their methodology has been controversial (cf. the Reinhart-Rogoff-Herndon-Ash debate), but the qualitative pattern survives subsequent reassessments and underpins the macro controls in our credit regression.

\textcite{borio2014bispaper} provides the BIS view of the same pattern: financial-cycle build-ups (credit-to-GDP gaps, asset-price gaps) precede sovereign-stress episodes by years, not months. The slow-moving nature of the cycle is the empirical regularity that motivates our prior that development-KPI signals price into credit before they price into equity.

The Bayesian sovereign-credit modelling literature is comparatively sparse; the bulk of credit-Bayesian work has focused on corporate bonds (\cite{bhojrajsengupta2003} on governance-rating links). Our hierarchical ordered-logistic specification in Section~\ref{sec:credit_method} is the sovereign analogue, with country-varying loadings under a partially-pooled prior.

\subsection{Local projections and panel causality}

\textcite{jorda2005} introduced local projections as a reduced-form alternative to vector autoregressions. Each horizon $h$ is a separate regression of the cumulative response on the shock, which sidesteps the recursion bias of VAR-implied IRFs --- the misspecification at horizon 1 propagates into all subsequent horizons in a VAR but is contained within a single regression in an LP. \textcite{jordataylor2016} extend the framework to event-study questions, applying it to fiscal consolidation with the augmented inverse-propensity-weighted estimator that corrects for selection on observables. \textcite{jordaschulariktaylor2013} apply the same machinery to the post-crisis growth experience of advanced economies and show that the depth of the credit boom predicts the depth of the post-bust recession at horizons up to five years.

The standard-error machinery for LP estimation in heterogeneous panels is the contribution of \textcite{driscollkraay1998}, whose covariance-matrix estimator handles arbitrary cross-sectional dependence by aggregating residuals over the time dimension and applying a Bartlett kernel. We deploy this estimator throughout the LP regressions in Section~\ref{sec:partI}. \textcite{cameronetal2008wildbootstrap} provide the wild-cluster bootstrap that we use as a small-sample-valid alternative when the Driscoll-Kraay kernel returns a non-PSD variance matrix.

\textcite{hamilton2018filter} argues forcefully that the cyclical component of a macro time series should be extracted by projecting $y_{t+h}$ on its own past values rather than by HP-filtering. The HP filter has known end-of-sample and bandwidth issues that have generated much subsequent criticism; the Hamilton (2018) alternative is now standard in macro-finance applications. Our LP regressions use Hamilton-cyclical sector-score innovations as the shock variable.

\textcite{abadiediamondhainmueller2010} introduce synthetic control as an identification method for case studies with a single treated unit and a small donor pool. The method has since been used to estimate the causal effect of trade liberalisations, immigration shocks, currency adoptions and many other ``policy episodes'' that cannot be analysed under conventional panel-DID. \textcite{abadie2021} surveys the synthetic-control literature and discusses the conditions under which the method delivers credible counterfactuals. Our Section~\ref{sec:synthrwanda} application to Rwanda's Vision 2020 is a direct deployment of the original method, with the placebo-permutation inference of \textcite{abadiediamondhainmueller2010}.

\textcite{dumitrescuhurlin2012} formalise the panel Granger non-causality test that we use in Section~\ref{sec:granger}. The key insight is that, in a heterogeneous panel, averaging the country-level Wald statistics --- rather than imposing a common slope --- preserves cross-country variation while still delivering a panel-level test. The asymptotic distribution of the average $\bar W$ statistic, under the null, is normal once standardised; this gives a tractable test in the small-$T$ panels that are common in development economics.

\textcite{pesaran2007} provides the cross-sectionally-augmented panel unit-root test (CIPS) that we run on the composite series and the GDP per capita series. \textcite{westerlund2007} provides the error-correction-based panel cointegration test that is now standard in macro panels. We compute both as a robustness layer in the analysis pipeline; the headline of the paper does not depend on a stationarity claim, so we leave the details to the appendix.

\subsection{Cross-country factor investing and emerging-market premia}

\textcite{famafrench2015} and the broader factor-investing literature establish that systematic factors --- value, momentum, size, profitability, investment --- exist across asset classes and countries. \textcite{asnessmoskowitzpedersen2013} explicitly document that value and momentum factors travel across countries and asset classes, including in emerging-market equity. \textcite{asnessfrazzinipedersen2019} formalise ``quality minus junk'' as a coherent factor that mixes profitability, growth and safety, and find that quality stocks earn higher returns than junk stocks with lower volatility.

The country-quality angle, as opposed to the firm-quality angle, has a thinner empirical track record. The published results on tradable performance of country-quality factors built on macro / governance / environmental variables are mixed. Our Section~\ref{sec:partIII} backtest looks carefully at the cap-weighted-vs-equal-weighted handicap implicit in the standard MSCI EM benchmark; the negative finding we report is consistent with the broader factor-investing literature's caveat that slow-moving fundamentals are not equity-tradable at monthly horizons.

\textcite{ilmanen2011} surveys expected-return premia across asset classes and emphasises that slow-moving fundamentals (cyclically-adjusted P/E, demographics, governance) tend to predict returns at long horizons but generate noisy signals at short horizons. The slow-moving nature of development indicators places them firmly in the long-horizon bin.

\textcite{politisromano1994} provide the stationary bootstrap that we use for Sharpe-ratio confidence intervals. The block-bootstrap family has become standard in performance-evaluation contexts where the underlying return series has serial dependence; with a mean block length of 12 months our resampling preserves the year-on-year autocorrelation in EM equity returns.

The randomly-selected-portfolio benchmark is a non-parametric approach that goes back to \textcite{ennisetal2013}: rather than comparing a strategy to a cap-weighted index, one compares it to the distribution of all strategies that could have been constructed under the same constraints. This is a more demanding benchmark because the index itself can be a bad benchmark when the cap-weighting scheme is not the relevant one. The literature on this is sparse but growing, and our 5\,000-random-basket exercise in Section~\ref{sec:partIII} is a deliberate borrow.

\subsection{Development economics and KPI frameworks}

Finally, the choice of which indicators to aggregate into a "development composite" is a question with a long history. \textcite{sen1999} and \textcite{nussbaum2011} argue for capability-based development metrics that look beyond GDP, and that argument has been operationalised by the UNDP's Human Development Index and by the OECD's Better Life Index. \textcite{acemoglurobinson2012} push the institutional-quality angle, emphasising that inclusive political and economic institutions are the long-run causal driver of development outcomes; governance indicators (our security-stability sector) are the proxy for that mechanism.

A growing practitioner literature combines sectoral KPIs into a country score, including handling of missing data, weighting of indicators and dynamic dashboards. Our pipeline mirrors that recipe with the addition of PCA-based weighting and the per-capita / logarithmic transforms.

The literature gap that this paper attempts to fill is in the joint treatment of (i) development KPIs as the conditioning information set, (ii) modern panel-causal methods for identification, (iii) Bayesian credit modelling, and (iv) tradable EM-factor backtesting. The four blocks have each been done well, but rarely together; the contribution of this work is in the cross.
